\section{Method}
\label{sec:method}

The method has four parts: a small vision-transformer encoder
(Section~\ref{sec:method_arch}), the Barlow Twins objective on two augmented
views (Sections~\ref{sec:method_loss} and~\ref{sec:method_views}), and a
training schedule with label-free checkpoint selection
(Section~\ref{sec:method_schedule}, Figure~\ref{fig:schedule}). The encoder
and the objective follow their published versions; the schedule is assembled
from standard, cited components. Deviations from published defaults are
disclosed where they occur.

\subsection{Encoder and projector}
\label{sec:method_arch}

The encoder is a ViT-Tiny~\cite{dosovitskiy2021vit,touvron2021deit} with
$8\times8$ patches operating on $64\times64$ images: 12 transformer layers,
192-dimensional embeddings, 3 attention heads, 5.4M parameters. The patch
size matches the native resolution of the pre-training images, giving
$8\times8=64$ patch tokens plus a classification (CLS) token. The
representation evaluated downstream is the CLS token after the final layer
normalization. During pre-training only, a projector maps this 192-dimensional
embedding to the space where the loss is computed: three linear layers
($192\rightarrow1024\rightarrow1024\rightarrow1024$), the first two followed
by batch normalization and ReLU. The projector follows the Barlow Twins
design~\cite{zbontar2021barlow} scaled down from its original 8192 dimensions
--- at 8192 the projector alone would exceed the encoder in parameter count
--- and is discarded after pre-training.

\subsection{Barlow Twins objective}
\label{sec:method_loss}

Barlow Twins makes the projector outputs of two augmented views of the same
images agree, while making the individual output dimensions carry
non-redundant information. Both pressures act on one matrix: the
cross-correlation between the two views' outputs, computed across the batch.
With $z^A, z^B \in \mathbb{R}^{B \times D}$ the batch-normalized projector
outputs of the two views ($B$ the batch size, $D=1024$ the projector width),
the cross-correlation matrix $\mathcal{C} \in \mathbb{R}^{D \times D}$ has
entries
\begin{equation}
  \mathcal{C}_{ij} \;=\;
  \frac{\sum_{b} z^{A}_{b,i}\, z^{B}_{b,j}}
       {\sqrt{\sum_{b} \bigl(z^{A}_{b,i}\bigr)^{2}}\,
        \sqrt{\sum_{b} \bigl(z^{B}_{b,j}\bigr)^{2}}},
  \label{eq:crosscorr}
\end{equation}
where $b$ indexes the batch and $i,j$ index output dimensions. The loss
pushes $\mathcal{C}$ toward the identity matrix,
\begin{equation}
  \mathcal{L}_{\mathrm{BT}} \;=\;
  \sum_{i} \bigl(1 - \mathcal{C}_{ii}\bigr)^{2}
  \;+\; \lambda \sum_{i} \sum_{j \neq i} \mathcal{C}_{ij}^{2},
  \label{eq:btloss}
\end{equation}
where the first term enforces invariance across views and the second
decorrelates output dimensions, with $\lambda$ weighting the two. A constant
representation cannot satisfy Equation~\ref{eq:btloss}: the diagonal of
$\mathcal{C}$ requires unit per-dimension correlation across views, so the
objective itself penalizes complete collapse, without negative pairs or
asymmetric networks~\cite{zbontar2021barlow}. We use the reference
implementation of the loss from the \texttt{lightly}
library~\cite{susmelj2020lightly} rather than reimplementing it, because the
loss has known silent failure modes (normalization along the wrong axis,
incorrect off-diagonal scaling) that a reference implementation rules out. We
set $\lambda = 10^{-2}$, deviating from the reference default $5 \cdot
10^{-3}$ that was tuned at projector width 8192~\cite{zbontar2021barlow}; our
value was chosen a priori alongside the narrower projector and was not
ablated.

\subsection{View generation}
\label{sec:method_views}

The two views are produced by the published two-view BYOL augmentation
pipeline~\cite{grill2020byol}, as implemented in
\texttt{lightly}~\cite{susmelj2020lightly}: random resized crop, horizontal
flip, color jitter (strength 1.0), and random grayscale; the second view
additionally solarizes with probability 0.2. We deviate from the published
recipe in two disclosed ways, both adaptations to $64\times64$ inputs. First,
Gaussian blur is disabled on both views: its kernel is sized for
$224\times224$ inputs. Second, the random crop retains at least 25\% of the
image area rather than the published 8\%: the crop scale is an area fraction,
so at $64\times64$ an 8\%-area crop keeps roughly $18\times18$ pixels ---
about $2\times2$ of our 8-pixel patches --- leaving two views little shared
content to align, and views that share too little information remove
task-relevant signal~\cite{tian2020infomin}. At $224\times224$, where the
recipe was tuned, the same fraction still keeps roughly $63\times63$ pixels.
Both bounds were set a priori and not ablated.

\subsection{Training schedule and checkpoint selection}
\label{sec:method_schedule}

Training follows a warmup-stable-decay (WSD)
schedule~\cite{hagele2024scaling,hu2024minicpm,zhai2022scalingvit}
(Figure~\ref{fig:schedule}). The learning rate warms up linearly over the
first 500 optimizer steps, then stays constant for the remainder of a fixed
epoch budget. We warm up in optimizer steps rather than the conventional ``10
epochs''~\cite{zbontar2021barlow,caron2021dino} because warmup stabilizes
optimizer updates, and an epoch-based rule degenerates across our set sizes
(at batch 250 it would give 40 warmup steps on 1k images and 4\,000 on 100k).

A constant rate, rather than the cosine decay common in self-supervised
pre-training, serves the selection mechanism: under a cosine schedule a
checkpoint read out well before the horizon attains higher loss than a model
whose schedule was matched to that
duration~\cite{hoffmann2022chinchilla,loshchilov2017sgdr}, so checkpoints
from different training stages would not be comparable, while under a
constant rate every checkpoint is schedule-equivalent. It also leaves runs
extendable, since no horizon is baked into the schedule. WSD schedules match
the final loss of a well-tuned cosine schedule~\cite{hagele2024scaling}, so
this convenience is not bought with accuracy --- though that evidence is from
language-model pre-training and its transfer to our setting is an assumption
we disclose (Section~\ref{sec:limitations}).

During the stable phase, a $k$-nearest-neighbour probe on held-out in-domain
data scores the encoder at a fixed epoch interval
(Section~\ref{sec:setup}), and the checkpoint with the best \emph{smoothed}
probe accuracy (mean of the last three probes) is retained. Smoothing
matters because selecting the maximum of a noisy, repeatedly evaluated signal
is biased toward flukes~\cite{cawley2010overfitting}. Selection is in-domain
--- the probe never sees the downstream dataset --- following standard
practice~\cite{caron2021dino}.

After the budget completes, a cooldown branch restarts from the selected best
checkpoint and anneals the learning rate to zero over $0.2 \times
t_{\mathrm{best}}$ steps (where $t_{\mathrm{best}}$ is the selected
checkpoint's step; floor 500 steps), with the $(1-\sqrt{t})$ decay shape that
performed best in the cooldown-shape ablation of H\"agele et
al.~\cite{hagele2024scaling}; decays of roughly 10\% of training suffice
where 2.5\% falls short~\cite{hu2024minicpm}. Branching the cooldown from a
checkpoint follows MiniCPM~\cite{hu2024minicpm}; rewinding to the
\emph{best-validation} checkpoint rather than the last one is our adaptation
(disclosed, not claimed standard): on small sets the probe accuracy can peak
and then decline within the budget, and annealing the last checkpoint would
anneal a degraded model. The annealed model is the one evaluated downstream.

Throughout, we train with AdamW~\cite{loshchilov2019adamw} (weight decay
$10^{-4}$, gradient-norm clipping 1.0) at batch size 250. AdamW replaces the
original recipe's LARS optimizer, which targets large-batch convolutional
training (batch 1k--32k); at our batch size, ViT self-supervised methods use
AdamW throughout~\cite{caron2021dino,chen2021mocov3,he2022mae}. The batch
size is 250 rather than the conventional 256 because 250 divides every
pre-training set size, so with a drop-last loader one epoch is exactly one
exposure of every image at every $N$ --- with batch 256, an epoch on the 1k
split would silently skip almost a quarter of it. The fixed
epoch budget is generous rather than calibrated: we do \emph{not} claim the
selected checkpoints are trained to convergence, and we disclose per set size
whether the validation signal plateaued within the budget
(Section~\ref{sec:results_curve}).

\begin{figure}[t]
  \centering
  % [PLACEHOLDER --- replaced by figures/schedule.pdf in the figure pass]
  \fbox{\parbox[c][3.2cm][c]{0.92\linewidth}{\centering\small
    Schedule schematic (pending): linear warmup $\rightarrow$ constant LR
    with periodic kNN probes $\rightarrow$ best smoothed-probe checkpoint
    marked $\rightarrow$ $(1-\sqrt{t})$ cooldown branched from it.}}
  \caption{The training schedule. The learning rate warms up for 500 steps
  and stays constant while a kNN probe periodically scores the encoder; after
  the fixed budget ends, training rewinds to the best-probe checkpoint and a
  short $(1-\sqrt{t})$ cooldown anneals the learning rate to zero. Every
  stable-phase checkpoint is schedule-comparable, and no training horizon is
  committed in advance.}
  \label{fig:schedule}
\end{figure}
